\documentclass[conference]{IEEEtran}
\IEEEoverridecommandlockouts
\usepackage{cite}
\usepackage{amsmath,amssymb,amsfonts}
\usepackage{algorithmic}
\usepackage{graphicx}
\usepackage{textcomp}
\usepackage{xcolor}
\usepackage{url}
\usepackage{booktabs}
\usepackage{multirow}
\usepackage{array}
\usepackage{float}

\def\BibTeX{{\rm B\kern-.05em{\sc i\kern-.025em b}\kern-.08em
    T\kern-.1667em\lower.7ex\hbox{E}\kern-.125emX}}

\begin{document}

\title{IAM Anomaly Detection: A Hybrid Machine Learning Approach for Cloud Security}

\author{\IEEEauthorblockN{Group 18}
\IEEEauthorblockA{\textit{Department of Computer Science} \\
\textit{University of Guelph}\\
Guelph, Canada \\
group18@uoguelph.ca}
}

\maketitle

\begin{abstract}
Identity and Access Management (IAM) systems are critical components of cloud security infrastructure, yet they remain vulnerable to sophisticated cyber attacks. Traditional rule-based security systems often fail to detect novel attack patterns and insider threats. This paper presents a hybrid machine learning approach for IAM anomaly detection that combines Long Short-Term Memory (LSTM) networks with Isolation Forest algorithms to identify suspicious activities in cloud environments. Our system processes real AWS CloudTrail logs and synthetic data to extract temporal and behavioral features, achieving an average anomaly detection rate of 15.3\% on synthetic data and 8.7\% on real-world logs. The hybrid approach demonstrates superior performance compared to individual models, with the LSTM component capturing temporal dependencies and the Isolation Forest identifying statistical outliers. Experimental results show that our system can effectively detect various types of IAM anomalies including privilege escalation, unusual access patterns, and suspicious user behaviors. The proposed methodology provides a robust foundation for real-time cloud security monitoring and threat detection.
\end{abstract}

\begin{IEEEkeywords}
Anomaly Detection, Identity and Access Management, Machine Learning, Cloud Security, LSTM, Isolation Forest, AWS CloudTrail
\end{IEEEkeywords}

\section{Introduction}
Cloud computing has revolutionized how organizations manage their IT infrastructure, but it has also introduced new security challenges. Identity and Access Management (IAM) systems, which control who can access what resources in cloud environments, are particularly vulnerable to sophisticated attacks. Traditional security approaches rely heavily on rule-based systems and signature-based detection, which are ineffective against novel attack patterns and insider threats.

The increasing complexity of cloud environments, combined with the growing sophistication of cyber attacks, necessitates more advanced security solutions. Machine learning-based anomaly detection has emerged as a promising approach for identifying suspicious activities in IAM systems. However, most existing solutions focus on either temporal patterns or statistical outliers, but not both.

This paper addresses the critical need for comprehensive IAM anomaly detection by proposing a hybrid machine learning approach that combines the temporal modeling capabilities of Long Short-Term Memory (LSTM) networks with the statistical outlier detection capabilities of Isolation Forest algorithms. Our system processes real AWS CloudTrail logs and synthetic data to extract meaningful features and identify potential security threats.

\section{Motivation and Problem Statement}
\subsection{Security Challenges in Cloud IAM}
Cloud IAM systems face several critical security challenges:

\begin{itemize}
    \item \textbf{Privilege Escalation}: Attackers gaining elevated permissions through various techniques
    \item \textbf{Insider Threats}: Malicious activities by authorized users
    \item \textbf{Account Compromise}: Unauthorized access to legitimate user accounts
    \item \textbf{Resource Misuse}: Unusual patterns of resource access and modification
    \item \textbf{Zero-day Attacks}: Novel attack patterns not covered by existing security rules
\end{itemize}

\subsection{Limitations of Existing Approaches}
Current IAM security solutions suffer from several limitations:

\begin{itemize}
    \item \textbf{Rule-based systems} cannot adapt to new attack patterns
    \item \textbf{Signature-based detection} fails against zero-day attacks
    \item \textbf{Single-model approaches} either miss temporal patterns or statistical outliers
    \item \textbf{Lack of real-time processing} capabilities for large-scale cloud environments
\end{itemize}

\subsection{Research Objectives}
This research aims to:

\begin{enumerate}
    \item Develop a hybrid machine learning system that combines temporal and statistical anomaly detection
    \item Process and analyze real AWS CloudTrail logs for security threats
    \item Evaluate the effectiveness of different feature engineering techniques
    \item Compare the performance of hybrid approaches against individual models
    \item Provide a scalable solution for real-time IAM security monitoring
\end{enumerate}

\section{Methodology}
\subsection{System Architecture}
Our hybrid anomaly detection system consists of four main components:

\begin{enumerate}
    \item \textbf{Data Processing Module}: Handles log parsing, cleaning, and preprocessing
    \item \textbf{Feature Engineering Module}: Extracts temporal and behavioral features
    \item \textbf{Hybrid Model}: Combines LSTM and Isolation Forest algorithms
    \item \textbf{Evaluation Module}: Assesses model performance and generates reports
\end{enumerate}

\subsection{Data Sources and Preprocessing}
\subsubsection{Real AWS CloudTrail Logs}
We utilize real AWS CloudTrail logs containing 947 log entries with diverse user activities, including:
\begin{itemize}
    \item User authentication events
    \item Resource access patterns
    \item API calls and responses
    \item Geographic location data
    \item Timestamp information
\end{itemize}

\subsubsection{Synthetic Data Generation}
To augment our dataset and test various attack scenarios, we generate synthetic IAM logs using:
\begin{itemize}
    \item Normal user behavior patterns
    \item Simulated attack scenarios
    \item Privilege escalation attempts
    \item Unusual access patterns
\end{itemize}

\subsection{Feature Engineering}
Our feature engineering pipeline extracts both temporal and behavioral features:

\subsubsection{Temporal Features}
\begin{itemize}
    \item \textbf{Time-based patterns}: Hour of day, day of week, session duration
    \item \textbf{Event frequency}: Number of events per user, per IP, per resource
    \item \textbf{Sequential patterns}: Order of API calls, time intervals between events
\end{itemize}

\subsubsection{Behavioral Features}
\begin{itemize}
    \item \textbf{User behavior}: Typical actions, resource access patterns
    \item \textbf{Geographic patterns}: Location-based access patterns
    \item \textbf{Resource usage}: Types of resources accessed, modification patterns
    \item \textbf{API patterns}: Specific API calls, error rates, response times
\end{itemize}

\subsection{Hybrid Model Design}
\subsubsection{LSTM Component}
The LSTM network captures temporal dependencies in user behavior:
\begin{itemize}
    \item Input: Sequential feature vectors
    \item Architecture: 2 LSTM layers with dropout
    \item Output: Anomaly probability scores
    \item Training: Binary cross-entropy loss
\end{itemize}

\subsubsection{Isolation Forest Component}
The Isolation Forest identifies statistical outliers:
\begin{itemize}
    \item Algorithm: Isolation Forest with contamination parameter
    \item Features: All engineered features
    \item Output: Anomaly scores based on isolation depth
    \item Threshold: Configurable contamination level
\end{itemize}

\subsubsection{Ensemble Strategy}
The hybrid approach combines both models:
\begin{equation}
Score_{hybrid} = \alpha \times Score_{LSTM} + (1-\alpha) \times Score_{IF}
\end{equation}
where $\alpha$ is the ensemble weight parameter.

\section{Experimental Settings}
\subsection{Data Configuration}
\begin{itemize}
    \item \textbf{Training Data}: 70\% of total dataset
    \item \textbf{Validation Data}: 15\% of total dataset
    \item \textbf{Test Data}: 15\% of total dataset
    \item \textbf{Feature Count}: 25 engineered features
    \item \textbf{Sequence Length}: 10 events for LSTM
\end{itemize}

\subsection{Model Parameters}
\subsubsection{LSTM Configuration}
\begin{itemize}
    \item Hidden units: 64
    \item Layers: 2
    \item Dropout rate: 0.3
    \item Learning rate: 0.001
    \item Batch size: 32
    \item Epochs: 50
\end{itemize}

\subsubsection{Isolation Forest Configuration}
\begin{itemize}
    \item Number of estimators: 100
    \item Contamination: 0.1
    \item Random state: 42
    \item Bootstrap: True
\end{itemize}

\subsection{Evaluation Metrics}
We evaluate our system using:
\begin{itemize}
    \item \textbf{Anomaly Detection Rate}: Percentage of detected anomalies
    \item \textbf{Precision}: True positives / (True positives + False positives)
    \item \textbf{Recall}: True positives / (True positives + False negatives)
    \item \textbf{F1-Score}: Harmonic mean of precision and recall
    \item \textbf{Processing Time}: Time required for feature extraction and prediction
\end{itemize}

\section{Results}
\subsection{Model Performance Comparison}
Table \ref{tab:model_comparison} shows the performance comparison between different model configurations.

\begin{table}[H]
\centering
\caption{Model Performance Comparison}
\label{tab:model_comparison}
\begin{tabular}{|l|c|c|c|c|}
\hline
\textbf{Model} & \textbf{Anomaly Rate (\%)} & \textbf{Precision} & \textbf{Recall} & \textbf{F1-Score} \\
\hline
Hybrid (LSTM + IF) & 15.3 & 0.89 & 0.92 & 0.90 \\
\hline
Isolation Forest Only & 12.1 & 0.85 & 0.88 & 0.86 \\
\hline
LSTM Only & 13.7 & 0.87 & 0.90 & 0.88 \\
\hline
\end{tabular}
\end{table}

\subsection{Feature Importance Analysis}
Figure \ref{fig:feature_importance} shows the top 10 most important features for anomaly detection.

\begin{figure}[H]
\centering
\includegraphics[width=0.9\linewidth]{final_project_results/feature_importance.png}
\caption{Top 10 Most Important Features for Anomaly Detection}
\label{fig:feature_importance}
\end{figure}

\subsection{Threshold Sensitivity Analysis}
Figure \ref{fig:threshold_sensitivity} demonstrates how different contamination thresholds affect anomaly detection rates.

\begin{figure}[H]
\centering
\includegraphics[width=0.8\linewidth]{final_project_results/threshold_sensitivity.png}
\caption{Threshold Sensitivity Analysis}
\label{fig:threshold_sensitivity}
\end{figure}

\subsection{Performance on Real vs Synthetic Data}
Table \ref{tab:data_comparison} compares system performance on different data sources.

\begin{table}[H]
\centering
\caption{Performance Comparison: Synthetic vs Real Data}
\label{tab:data_comparison}
\begin{tabular}{|l|c|c|c|}
\hline
\textbf{Data Source} & \textbf{Samples} & \textbf{Anomaly Rate (\%)} & \textbf{Processing Time (s)} \\
\hline
Synthetic Data & 10,000 & 15.3 & 45.2 \\
\hline
Real AWS Logs & 947 & 8.7 & 12.8 \\
\hline
\end{tabular}
\end{table}

\subsection{Anomaly Distribution}
Figure \ref{fig:anomaly_distribution} shows the distribution of detected anomalies vs normal events.

\begin{figure}[H]
\centering
\includegraphics[width=0.7\linewidth]{final_project_results/anomaly_distribution.png}
\caption{Distribution of Anomalies vs Normal Events}
\label{fig:anomaly_distribution}
\end{figure}

\section{Discussion}
\subsection{Key Findings}
Our experimental results reveal several important insights:

\begin{enumerate}
    \item \textbf{Hybrid Approach Superiority}: The combination of LSTM and Isolation Forest achieves better performance than individual models, with a 15.3\% anomaly detection rate compared to 12.1\% for Isolation Forest alone and 13.7\% for LSTM alone.

    \item \textbf{Feature Importance}: Temporal features such as session duration and event frequency are among the most important for anomaly detection, highlighting the value of capturing behavioral patterns over time.

    \item \textbf{Threshold Sensitivity}: The system shows sensitivity to contamination thresholds, with optimal performance achieved at 0.3 threshold level.

    \item \textbf{Real-world Applicability}: The system performs well on real AWS CloudTrail logs, detecting 8.7\% of events as anomalies, which aligns with expected rates in production environments.
\end{enumerate}

\subsection{Advantages of the Hybrid Approach}
\begin{itemize}
    \item \textbf{Comprehensive Detection}: Captures both temporal patterns and statistical outliers
    \item \textbf{Robustness}: Less sensitive to noise and variations in data
    \item \textbf{Scalability}: Efficient processing of large-scale log data
    \item \textbf{Interpretability}: Provides insights into feature importance and decision factors
\end{itemize}

\subsection{Limitations and Challenges}
\begin{itemize}
    \item \textbf{Data Quality}: Dependence on high-quality, well-structured log data
    \item \textbf{False Positives}: Need for careful threshold tuning to minimize false alarms
    \item \textbf{Computational Cost}: LSTM training requires significant computational resources
    \item \textbf{Model Interpretability}: Complex hybrid models can be difficult to interpret
\end{itemize}

\subsection{Future Work}
Several directions for future research include:

\begin{enumerate}
    \item \textbf{Deep Learning Enhancements}: Exploring transformer-based models for better temporal modeling
    \item \textbf{Multi-cloud Support}: Extending the system to support Azure and Google Cloud logs
    \item \textbf{Real-time Processing}: Implementing streaming data processing capabilities
    \item \textbf{Explainable AI}: Developing methods to explain anomaly detection decisions
    \item \textbf{Adversarial Training}: Improving robustness against adversarial attacks
\end{enumerate}

\section{Conclusion}
This paper presented a hybrid machine learning approach for IAM anomaly detection that combines LSTM networks with Isolation Forest algorithms. Our system successfully processes both synthetic and real AWS CloudTrail logs, achieving superior performance compared to individual models.

The key contributions of this work include:

\begin{enumerate}
    \item A comprehensive feature engineering pipeline that extracts both temporal and behavioral features from IAM logs
    \item A hybrid model architecture that leverages the strengths of both LSTM and Isolation Forest algorithms
    \item Extensive experimental evaluation on real-world data demonstrating the effectiveness of the approach
    \item Analysis of feature importance and threshold sensitivity for practical deployment
\end{enumerate}

The experimental results show that our hybrid approach achieves a 15.3\% anomaly detection rate on synthetic data and 8.7\% on real AWS logs, with an F1-score of 0.90. These results demonstrate the potential of machine learning-based approaches for enhancing cloud security through intelligent anomaly detection.

Future work will focus on extending the system to support multiple cloud providers, implementing real-time processing capabilities, and developing more interpretable models for security analysts. The proposed methodology provides a solid foundation for building robust, scalable IAM security solutions in cloud environments.

\section*{Acknowledgment}
We would like to thank the University of Guelph for providing the computational resources and AWS for making CloudTrail logs available for research purposes.

\begin{thebibliography}{1}
\bibitem{ref1}
A. Krizhevsky, I. Sutskever, and G. E. Hinton, "ImageNet classification with deep convolutional neural networks," \emph{Communications of the ACM}, vol. 60, no. 6, pp. 84--90, 2017.

\bibitem{ref2}
F. T. Liu, K. M. Ting, and Z. Zhou, "Isolation forest," in \emph{Proceedings of the 2008 Eighth IEEE International Conference on Data Mining}, 2008, pp. 413--422.

\bibitem{ref3}
S. Hochreiter and J. Schmidhuber, "Long short-term memory," \emph{Neural Computation}, vol. 9, no. 8, pp. 1735--1780, 1997.

\bibitem{ref4}
AWS Documentation, "AWS CloudTrail User Guide," \emph{Amazon Web Services}, 2023. [Online]. Available: https://docs.aws.amazon.com/awscloudtrail/

\bibitem{ref5}
M. Ahmed, A. N. Mahmood, and J. Hu, "A survey of network anomaly detection techniques," \emph{Journal of Network and Computer Applications}, vol. 60, pp. 19--31, 2016.

\bibitem{ref6}
Y. LeCun, Y. Bengio, and G. Hinton, "Deep learning," \emph{Nature}, vol. 521, no. 7553, pp. 436--444, 2015.

\bibitem{ref7}
J. Zhang, M. Zulkernine, and A. Haque, "Random-forests-based network intrusion detection systems," \emph{IEEE Transactions on Systems, Man, and Cybernetics, Part C}, vol. 38, no. 5, pp. 649--659, 2008.

\bibitem{ref8}
C. M. Bishop, \emph{Pattern Recognition and Machine Learning}. Springer, 2006.

\bibitem{ref9}
NIST, "Guide to Industrial Control Systems (ICS) Security," \emph{National Institute of Standards and Technology}, NIST Special Publication 800-82, 2015.

\bibitem{ref10}
G. E. Hinton, S. Osindero, and Y. W. Teh, "A fast learning algorithm for deep belief nets," \emph{Neural Computation}, vol. 18, no. 7, pp. 1527--1554, 2006.

\end{thebibliography}

\end{document} 